\subsection{Permutations, symmetries and cycles}
A permutation, is a bijective function from a set to itself. The set of all permutations of 
a fine set $A$ is a group under the operation of composition. We will see how this group is 
an extremely important object in Galois theory, as it is used to study the symmetries of the 
roots of polynomial equations. 

\begin{definition}[Permutation]
A permutation of a set $A$ is a bijective function $\sigma: A \to A$. 
\end{definition} \vspace{0.2em}

\begin{definition}[Symmetric group]
The set $Sn$ is the set of all the permutations of the set $\{1,2,\dots,n\}$. Please 
notice that the such set is indeed isomorphic to the set of all permutations of the 
set $A$, for every choice of $A$ such that it has finite cardinality $n$.
\end{definition} \vspace{0.4em}

We can indicate a permutation $\sigma$ in $Sn$ with the following notation (matrix notation):

$$ \sigma = \begin{pmatrix} 
    1 & 2 & 3 & \dots & n \\ 
    \sigma(1) & \sigma(2) & \sigma(3) & \dots & \sigma(n) 
\end{pmatrix} $$ \vspace{0.2em}

\begin{definition}[Cycle]
A cycle of length $k$ is a permutation $\sigma$ such that 
every element gets mapped onto itself after $k$ mappings (Repeatedly applying
the permutation $\sigma$ $k$ times, can be written as $ \sigma^{k}$ since 
in a group the operation is typically see as multiplication-alike):

$$ \sigma^{k}(n) = n $$ \vspace{0.2em}
\end{definition} \vspace{0.2em}

We can indicate a cycle $\tau$ of length $k$ with the following notation (tuple notation):

$$ \tau = (a_0,\; a_1,\; \dots,\; a_k ) $$ \vspace{0.2em}

Which represents the permutation described by the following case-by-case definition:

$$ \tau(x) = \begin{cases}
        a_{i+1} & \text{if } x = a_i \land x \neq k \\
        a_1 & \text{if } i = k \\
        x & \text{otherwise } \\
    \end{cases} 
$$ \vspace{0.2em}

\newpage

Notice that $(a_0)$ can be written as $(a_0,\; a_0)$, and that's the identity in every $S_n$. \\

\begin{lemma} 
    $S_n$ forms the group $(S_n, \circ)$ under the operation of composition.
\end{lemma} 
\begin{proof}
    It's clearly closed, since the product of two permutations is a
    permutation. The identity permutation is the one that maps every element to itself,
    and since every permutation is a bijection, it has an inverse. Also, by the definition of
    composition, the operation is associative. Thus, $S_n$ is a group under the operation of
    composition.
\end{proof} \vspace{0.2em}

\begin{lemma}
\label{PermutationsCancellationLaw}
    Given $S_n \subseteq S_{n+1}$, we can express every permutation in $S_{n+1}$ as the 
    product of a permutation in $S_n$ and a length-2-cycle in the form $(n+1,\;\; 1 \leq i \leq n)$.
\end{lemma}  
\begin{proof}
    Let's denote the length-2-cycle as $\beta_{i} = (n+1,\; i)$. There exist then exactly $n$ 
    different functions $f_i(\sigma) = \sigma \circ \beta_{i}$, where $i$ is less then or equal 
    to $n$. \\

    Now, observe that the function $f$ is injective, since the cancellation law holds for
    the group of permutations just like any other group (lemma \ref{CancellationLaw}).

    $$ f_i(\sigma) = f_i(\tau) \Rightarrow \sigma \circ \beta = \tau \circ \beta $$
    $$ \Rightarrow \sigma \circ \beta \circ \beta^{-1} = \tau \circ \beta \circ \beta^{-1} $$
    $$ \Rightarrow \sigma = \tau $$ \vspace{0.2em}

    Also notice that the images of the function $f_i$ are all distinct since they all 
    maps to element that have been multiplied by a different factor which is not in the domain. Thus: \\
    
    $$ |\bigcup_{1 \leq i \leq n} \vv{f_i}(S_i) | = \sum_{1 \leq i \leq n} |\vv{f_i}(S_n)| 
    = (n + 1) \cdot n! = (n+1)! $$ \vspace{0.2em}

    For finte sets, we can say that $H \subseteq G$ is equal to $G$ when $|H| = |G|$, Hence the union 
    of all the products of permutations in $S_n$ with a cycle in $S_{n+1}$ is itself just $S_{n+1}$.
\end{proof} \vspace{0.2em}

\begin{theorem}[Permutations-Factorization] 
\label{PermutationsFactorization}
    Every permutation can be expressed as a product of cycles 
    of length $2$, which we call prime-factors of the permutation.
\end{theorem} 
\begin{proof}
    Let's proceed by 
    induction on the number of elements in the set $A$. The base case is when $A$ has only 
    two elements is trivially true, since the only permutations are the $(1,2)$ cycle and 
    the identity, which can be expressed as the product $(1,2) \circ (2,1)$. \\

    Now, for every permutation in $Sn$ assume that a factorization exists. Observe that 
    that every permutation in $Sn$ can be expressed as a product of length-cycles, and such product 
    of cycles also describes a permutation in $S_{n+1}$. \\
    
    Since by the result of lemma \ref{PermutationsCancellationLaw} we show that 
    $S_{n+1}$ can itself be described as the group of all the products of permutations in 
    $S_{n}$ with a length-2-cycle, every permutation in $S_{n+1}$ is herby the product of 
    a sequence of length-2-cycles. \\
\end{proof} \vspace{0.2em}

\newpage